\documentclass[10pt,a4paper]{amsart}
\usepackage[margin=19mm]{geometry}
\usepackage{amsmath,amssymb,mathtools}
\usepackage[scr=rsfs,cal=euler]{mathalfa}
\usepackage{xfrac}
\usepackage[unicode,hidelinks,bookmarks=false]{hyperref}
\newcommand{\ie}{i.e.\ }
\newcommand{\cf}{cf.\ }
\newcommand{\resp}{respectively\ }
\newcommand{\Subsection}{\subsection}
\providecommand{\nobreakdash}{\nobreak\hbox{-}}
\setlength{\emergencystretch}{2em}
\multlinegap0pt
% ---------------------------------------------------------------------------
% Editorial AMS wrapper prelude. The theorem environments and the mathematical macros below are
% transcribed from the paper's own preamble (cached-originals/source0.tex lines 1--64); the AMS amsart
% class, the named format packages of the pinned TeX Live runtime and this editorial formatting are not
% part of the author's text. No mathematics is changed.
% ---------------------------------------------------------------------------
\usepackage{needspace}
\allowdisplaybreaks[2]
\numberwithin{equation}{section}
\newtheorem{theorem}{Theorem}[section]
\newtheorem{proposition}[theorem]{Proposition}
\newtheorem{lemma}[theorem]{Lemma}
\newtheorem{corollary}[theorem]{Corollary}
\theoremstyle{definition}
\newtheorem{definition}[theorem]{Definition}
\newtheorem{assumption}[theorem]{Assumption}
\theoremstyle{remark}
\newtheorem{remark}[theorem]{Remark}
\newtheorem*{remark*}{Remark}
\newcommand{\R}{\mathbb R}
\newcommand{\T}{\mathbb T}
\newcommand{\cS}{\mathcal S}
\newcommand{\eps}{\varepsilon}
\newcommand{\etae}{\eta_{\eps}}
\newcommand{\de}{\delta_{\eps}}
\newcommand{\we}{w_{\eps}}
\newcommand{\tr}{\operatorname{tr}}
\newcommand{\dist}{\operatorname{dist}}
\newcommand{\loc}{\mathrm{loc}}
\newcommand{\comp}{\mathrm{comp}}
\newcommand{\tot}{\mathrm{tot}}
\newcommand{\C}{\mathbb C}
\newcommand{\dd}{\,dx}
\newcommand{\calA}{\mathcal A}
\DeclareMathOperator{\Res}{Res}
\DeclareMathOperator{\osc}{osc}

\begin{document}
\title{Non-simple blow-up for the Chern--Simons--Higgs equation: A priori analysis and constructions}
\author{Youngae Lee and Lei Zhang}
\date{Source: arXiv:2609.20042v1; distributed under CC BY 4.0}
\maketitle
\noindent\emph{Source, version and licence.} \emph{Non-simple blow-up for the Chern--Simons--Higgs equation: A priori analysis and constructions}, by Youngae Lee and Lei Zhang. The source of this editable unit is the e-print of \textbf{version v1} of arXiv:2609.20042, https://arxiv.org/abs/2609.20042v1, whose material is distributed under the Creative Commons Attribution 4.0 International licence, http://creativecommons.org/licenses/by/4.0/, as recorded in the arXiv licence metadata for this paper and shown on that abstract page. The whole of the body below is version v1, the newest version of the paper. Full rights notice: licenses/arxiv-2609.20042-CC-BY-4.0.txt.\par\smallskip
\noindent\textit{Editorial scope.} Counted unit: original Theorem 5.1 of the article, the statement carried by the source environment \texttt{theorem} with source label \texttt{prop:refined-profile}, cached-originals/source0.tex lines 1307--1394, printed as \textbf{Theorem 5.1} exactly as in the article: uniformly in the bubble index, the solution of the Chern-Simons-Higgs equation is approximated at each bubble point by the corresponding single-bubble profile together with the stated derivative estimates, in the refined form of the three assertions of the statement. It is delivered together with the authors' own complete proof as the single primary proof body, source0.tex lines 2852--3378 (527 lines): the \texttt{proof} environment that belongs to the statement and establishes it in full. No proof marker is added and no mathematics is written, repaired or re-derived; the authors' notation and the printed slips of the source are kept literal. Necessary complete context is retained uncounted, in source order: lines 538--547, the source \texttt{lemma} with source label \texttt{thm:linearized-classification} of the article that the retained passages refer to; lines 571--585, the source \texttt{lemma} with source label \texttt{lem:weighted-green} of the article that the retained passages refer to; lines 2654--2661, the source \texttt{lemma} with source label \texttt{lem:preliminary-center} of the article that the retained passages refer to; lines 2680--2689, the source \texttt{lemma} with source label \texttt{lem:preliminary-pointwise} of the article that the retained passages refer to; lines 2775--2784, the source \texttt{lemma} with source label \texttt{lem:first-moment} of the article that the retained passages refer to; lines 2796--2799, the source \texttt{equation} with source label \texttt{eq:correct-source-change-variable} of the article that the retained passages refer to; lines 2850--2850, the source \texttt{passage} with source label \texttt{app:pointwise-proof} of the article that the retained passages refer to. The article's own numbering is reproduced by explicit counter settings: the theorem-like counter of the article is set before each retained passage, so the counted statement prints with the number the article gives it. No \texttt{\textbackslash cite} command occurs in any retained passage: the closure of the retained passages over references and citations was computed, it contains no citation at all, so this unit delivers no bibliography entry and the article's own bibliography data is neither spliced into the bundled source nor redistributed. The source's own class is \texttt{amsart}; the e-print ships no class or style file, so no publisher class is shipped, used or redistributed, and the wrapper is the AMS \texttt{amsart} class with the article's own macros and theorem declarations transcribed from its preamble. The retained passages contain no figure environment and no graphics-inclusion command, so no artwork is redistributed. Every declared body of this unit is an exact, unmodified span of source0.tex: no body is a bounded passage comparison, \texttt{omit} was not used anywhere, and consequently no fidelity receipt is asserted in delivery-evidence.json. The complete concatenated original source is bundled as sources/210831/source0.tex. Omitted from this editable unit and therefore absent from it are source0.tex lines 1--537; 548--570; 586--1306; 1395--2653; 2662--2679; 2690--2774; 2785--2795; 2800--2849; 3379--3652. The AMS \texttt{amsart} wrapper, the named format packages of the pinned TeX Live runtime and this editorial source, licence and scope paragraph are editorial formatting only.\par\smallskip
\setcounter{section}{3}
\setcounter{equation}{4}
\setcounter{theorem}{0}
\begin{lemma}[Sublinear kernel and normalization]
\label{thm:linearized-classification}
If $L_\alpha\phi=0$ in $\R^2$ and
$|\phi(x)|\le C(1+|x|)^\tau$ for some $0<\tau<1$, then
\begin{equation}\label{eq:kernel-classification}
 \phi=c_0\partial_\alpha U_\alpha
       +c_1\partial_{x_1}U_\alpha+c_2\partial_{x_2}U_\alpha.
\end{equation}
In particular, $\phi(0)=0$ and $\nabla\phi(0)=0$ imply $\phi=0$.
\end{lemma}

\setcounter{section}{3}
\setcounter{equation}{7}
\setcounter{theorem}{1}
\begin{lemma}[Weighted Green estimate]\label{lem:weighted-green}
Let $G_R$ be the positive Dirichlet Green function of $-\Delta$ in
$B_R$, $R\ge2$. For $0<\sigma<1$, $\rho>0$, and every $x\in B_R$,
\begin{equation}\label{eq:weighted-green}
 \int_{B_R}|G_R(x,y)-G_R(0,y)|(1+|y|)^{-2-\rho}\,dy
 \le C_{\rho,\sigma}(1+|x|)^\sigma.
\end{equation}
Also,
\begin{equation}\label{eq:weighted-green-gradient}
 \int_{B_R}|\nabla_xG_R(x,y)|(1+|y|)^{-2-\rho}\,dy
 \le C_{\rho,\sigma}(1+|x|)^{\sigma-1}.
\end{equation}
The constants are independent of $R$; the same bounds hold for disks
whose centers are $o(1)$ from the normalization point.
\end{lemma}

\setcounter{section}{5}
\setcounter{equation}{0}
\setcounter{theorem}{0}
\begin{theorem}[Pointwise approximation at each Chern--Simons bubble]
\label{prop:refined-profile}
Assume the finite-height cluster setting and the local normalizations above.
Fix
\begin{equation}
  0<\sigma<\min\{1,m_*-4\}.
  \label{eq:sigma-choice}
\end{equation}
Uniformly for $t=0,\dots,S-1$, the following assertions hold.

\smallskip
\noindent\textup{(i)}   
For $|z|\le r/(2s_{t,\eps})$, we have
\begin{equation}\label{eq:global-pointwise-O2}
 |W_{t,\eps}(z)|+|\nabla W_{t,\eps}(z)|
 \le C\bigl(\etae^2+\etae\de^2\bigr).
\end{equation}
If $\varphi_i$ is constant on the local chart, the term
$\etae\de^2$ in \eqref{eq:global-pointwise-O2} is absent.

\smallskip
\noindent\textup{(ii)} \(
  \left|
  \nabla\log H_{t,\eps}(q_{t,\eps}^*)
  +\frac{m_{t,\eps}}4
   \de\nabla_y\varphi_i(y)\Big|_{y=\de q_{t,\eps}^*}
  \right|
  \le C\etae^2.
\)

\smallskip
\noindent\textup{(iii)}
Writing
$\mathbf{d}_{t,\eps}=q_{t,\eps}^*-q_{t,\eps}$ and
$\kappa_{t,\eps}=F(\alpha_{t,\eps})$, one has
\(
  \mathbf{d}_{t,\eps}
  =-\frac{2s_{t,\eps}^2}{\kappa_{t,\eps}}
  \nabla\mathfrak h_{t,\eps}(q_{t,\eps}^*)
  +O(\etae^3+\etae^4\de^2).
  \)
In particular, $|q_{t,\eps}^*-q_{t,\eps}|\le C\etae^2$.

\smallskip
\noindent\textup{(iv)}
Let
\begin{equation}\label{eq:B-tracefree}
\begin{aligned}
 \mathsf B_{t,\eps}
   &:=D^2\log H_{t,\eps}(q_{t,\eps}^*),
 &\mathsf B_{t,\eps}^{\circ}
   &:=\mathsf B_{t,\eps}-\tfrac12(\tr\mathsf B_{t,\eps})I_2,\\
 \mathsf C_{t,\eps}
   &:=\de^2D^2_y\varphi_i(y)\Big|_{y=\de q^*_{t,\eps}},
 &\mathsf C_{t,\eps}^{\circ}
   &:=\mathsf C_{t,\eps}-\tfrac12(\tr\mathsf C_{t,\eps})I_2.
\end{aligned}
\end{equation}
There is a unique bounded pure second Fourier mode $\Psi_{2,t,\eps}$
satisfying
\begin{equation}\label{eq:Psi2-equation}
\left\{
\begin{aligned}
 L_{\alpha_{t,\eps}}\Psi_{2,t,\eps}
   &=-\tfrac12(z^T\mathsf B_{t,\eps}^{\circ}z)F'(U_{t,\eps}(z))
     -\tfrac12(z^T\mathsf C_{t,\eps}^{\circ}z)F(U_{t,\eps}(z)),\\
 \Psi_{2,t,\eps}(0)&=0,\qquad \nabla\Psi_{2,t,\eps}(0)=0.
\end{aligned}
\right.
\end{equation}
With the conformal first Fourier correction $\Xi_{1,t,\eps}$
defined in \eqref{eq:Xi1-definition} of
Appendix~\ref{app:pointwise-proof},
for $|z|\le 3r/(4s_{t,\eps})$ we have
\begin{equation}\label{eq:second-mode-refinement}
 \left|\widehat v_{t,\eps}(z)-U_{t,\eps}(z)
       -\Xi_{1,t,\eps}(z)-s_{t,\eps}^2\Psi_{2,t,\eps}(z)\right|
 \le C\bigl(\etae^{2+\sigma}+\etae^2\de^2\bigr)(1+|z|)^\sigma.
\end{equation}

\smallskip
\noindent\textup{(v)} We have
\begin{equation}
  \widetilde m_{t,\eps}=m_{t,\eps}
  +O(\etae^{2+\sigma}+\etae^2\de^2).
  \label{eq:local-mass-refined}
\end{equation}
\end{theorem}

\setcounter{section}{8}
\setcounter{equation}{7}
\setcounter{theorem}{0}
\begin{lemma}[Preliminary comparison of the two centers]
\label{lem:preliminary-center}
With the notation above,
\begin{equation}
  |q_{t,\eps}^*-q_{t,\eps}|\le C\etae^2.
  \label{eq:preliminary-center}
\end{equation}
\end{lemma}

\setcounter{section}{8}
\setcounter{equation}{9}
\setcounter{theorem}{1}
\begin{lemma}[Preliminary weighted estimate]
\label{lem:preliminary-pointwise}
Under the assumptions of Theorem~\ref{prop:refined-profile},
\begin{equation}\label{eq:preliminary-pointwise}
 |W_{t,\eps}(z)|\le C\etae(1+|z|)^\sigma,
 \qquad |\nabla W_{t,\eps}(z)|\le C\etae
 \quad\left(|z|\le\frac{3r}{4s_{t,\eps}}\right).
\end{equation}
In fact the estimate for $W_{t,\eps}$ holds on the full rescaled disk.
\end{lemma}

\setcounter{section}{8}
\setcounter{equation}{15}
\setcounter{theorem}{2}
\begin{lemma}[First moment on a tangent circle]\label{lem:first-moment}
Suppress the core indices, including $\mathfrak u=\mathfrak u_{t,\eps}$.
Put $\mathbf{d}=q^*-q$, $r_\eps=r-|\mathbf{d}|$, and
$\boldsymbol\omega(\theta)=(\cos\theta,\sin\theta)$.
Then
\begin{equation}\label{eq:first-moment-coarse}
 \frac1{2\pi}\int_0^{2\pi}\mathfrak u(q^*+r_\eps \boldsymbol\omega(\theta))\boldsymbol\omega(\theta)\,d\theta
 =-\frac{m}{2r}\mathbf{d}+O(\etae^2).
\end{equation}
\end{lemma}

\setcounter{section}{8}
\setcounter{equation}{16}
\setcounter{theorem}{3}
\begin{equation}\label{eq:correct-source-change-variable}
 f_\eps(y)\,dy=\mathcal F_\eps(z)\,dz,\qquad
 \mathcal F_\eps:=e^\Gamma F(U+W+\ell).
\end{equation}

\setcounter{section}{8}
\setcounter{equation}{21}
\setcounter{theorem}{3}
\label{app:pointwise-proof}

\setcounter{section}{8}
\setcounter{equation}{21}
\setcounter{theorem}{3}
\begin{proof}[Proof of Theorem~\ref{prop:refined-profile}]
We suppress the core indices and put $\kappa=F(\alpha)$,
$\mathbf{d}=q^*-q$, and
{ 
\begin{equation}
  \mathbf{b}:=\nabla\log H(q^*),
  \qquad
  \mathbf{g}:=\de\nabla_y\varphi_i(y)\Big|_{y=\de q^*},
  \qquad
  \mathbf{b}^{\rm eff}:=\mathbf{b}+\frac m4\mathbf{g},
  \label{eq:b-g-effective}
\end{equation}
}
The conformal normalization gives
\begin{equation}
  \mathbf{g}=O(\de^2),
  \qquad
  s_\eps=\etae(1+O(\de^2)).
  \label{eq:s-g-size}
\end{equation}

Lemma~\ref{lem:preliminary-center} gives $\mathbf{d}=O(\etae^2)$ and
$\zeta_\eps=-\mathbf{d}/s_\eps=O(\etae)$. The maximum condition is
\begin{equation}
  0
  =
  s_\eps\nabla v_\eps(q)
  =
  \nabla\widehat v(\zeta_\eps)
  +s_\eps\nabla\mathfrak h(q).
  \label{eq:max-condition-scaled}
\end{equation}

\needspace{5\baselineskip}
\medskip
\noindent
{ 
\emph{Step 1: the conformal first Fourier correction.}
}

Taylor expansion gives
{ 
\begin{align}
  \log A(z)
  &=
  s_\eps \mathbf{b}\cdot z
  +O(s_\eps^2|z|^2),
  \ \
  \varphi_i(\de(q^*+s_\eps z))
  -\varphi_i(\de q^*)
  =
  s_\eps \mathbf{g}\cdot z
  +O(\de^2s_\eps^2|z|^2).
  \label{eq:conformal-first}
\end{align}
}
By Lemma~\ref{lem:preliminary-pointwise},
{ 
\begin{align}
  L_\alpha W
  ={}&
  -s_\eps(\mathbf{b}\cdot z)F'(U_\alpha)
  -s_\eps(\mathbf{g}\cdot z)F(U_\alpha)
  +E_\eps^{(2)},
  \label{eq:W-first-decomposition}
\end{align}
}
where \begin{equation}
\begin{aligned}
  |E_\eps^{(2)}(z)|
  &\le 
  c\left(\eta_\eps^2(1+\de^2)
  (1+|z|)^{2-m_*}
 +
  (1+|z|)^{-m_*}|W(z)|^2\right)  \le C \eta_\eps^2
  (1+|z|)^{2-m_*}.
  \label{eq:E2-crude}
\end{aligned}\end{equation}

For the coefficient first mode set
\begin{equation}
  Z_{\mathbf{b}}(z)
  :=
  -s_\eps \mathbf{b}\cdot z
  -\frac{2s_\eps}{\kappa}
   \mathbf{b}\cdot\nabla U_\alpha(z).
  \label{eq:explicit-first-mode}
\end{equation}
Then
\(L_\alpha Z_{\mathbf{b}} = -s_\eps(\mathbf{b}\cdot z)F'(U_\alpha).\)

{ 
For the conformal first mode, let $Y_\alpha$ be the normalized regular
solution of
\begin{equation}
  Y_\alpha''
  +\frac1{\varrho}Y_\alpha'
  -\frac1{\varrho^2}Y_\alpha
  +F'(U_\alpha)Y_\alpha
  =
  -\varrho F(U_\alpha),
  \qquad
  Y_\alpha(0)=Y_\alpha'(0)=0.
  \label{eq:Y-first-mode}
\end{equation}}
The normalized solution has the explicit form
\begin{equation}\label{eq:Y-explicit}
 Y_\alpha(\varrho)=\frac{\varrho^2}{4}U_\alpha'(\varrho).
\end{equation}
Indeed, let $L_{\alpha,1}=\partial_\varrho^2+
\varrho^{-1}\partial_\varrho-\varrho^{-2}+F'(U_\alpha)$.
Since $L_{\alpha,1}U_\alpha'=0$, the radial profile equation gives
\[
 L_{\alpha,1}(\varrho^2U_\alpha')
 =4(\varrho U_\alpha''+U_\alpha')
 =-4\varrho F(U_\alpha).
\]
Formula~\eqref{eq:Y-explicit} has the required normalization at zero;
uniqueness follows because the regular homogeneous first mode is a
multiple of $U_\alpha'$, whose derivative at zero is $-\kappa/2\ne0$.
It also gives, with differentiated expansions,
\[
 \begin{aligned}
 Y_\alpha(\varrho)&=-\kappa\varrho^3/8+O(\varrho^5)
 &&(\varrho\to0),\quad\
 Y_\alpha(\varrho)&=-\frac m4\varrho+O(\varrho^{3-m_*})
 &&(\varrho\to\infty).
 \end{aligned}
\]
Consequently
\begin{equation}\label{eq:Q-first-mode}
 \begin{aligned}
 Q_\alpha(\varrho)
 &:=Y_\alpha(\varrho)+\frac m4\varrho
       +\frac{m}{2\kappa}U_\alpha'(\varrho) =\left(\frac{\varrho^2}{4}+\frac{m}{2\kappa}\right)
       U_\alpha'(\varrho)+\frac m4\varrho
 \end{aligned}
\end{equation}
satisfies $Q_\alpha=O(\varrho^3)$, $Q_\alpha'=O(\varrho^2)$ at zero
and $Q_\alpha=O(\varrho^{-1})$, $Q_\alpha'=O(\varrho^{-2})$ at infinity.

Set
{ 
\begin{equation}
  \Xi_1(z)
  :=
  s_\eps Q_\alpha(|z|)
  \mathbf{g}\cdot\frac z{|z|}.
  \label{eq:Xi1-definition}
\end{equation}
}
Then
\begin{equation}
  |\Xi_1(z)|+|\nabla\Xi_1(z)|
  \le C\etae\de^2.
  \label{eq:Xi1-bound}
\end{equation}

With $Z_{\mathbf g}(z):=s_\eps Y_\alpha(|z|)\mathbf g\cdot z/|z|$,
the definitions of $\mathbf b^{\rm eff}$ and $Q_\alpha$ give
\begin{align}
  Z_{\mathbf{b}}(z)+Z_{\mathbf{g}}(z)
  ={}&
  -s_\eps \mathbf{b}^{\rm eff}\cdot z
  -\frac{2s_\eps}{\kappa}
   \mathbf{b}^{\rm eff}\cdot\nabla U_\alpha(z)
  +\Xi_1(z).
  \label{eq:first-mode-effective-decomposition}
\end{align}

Let \(\widetilde W:=W-\Xi_1\).
Then \eqref{eq:W-first-decomposition} implies that
{ 
\begin{equation}
  L_\alpha\widetilde W
  =
  -s_\eps(\mathbf{b}^{\rm eff}\cdot z)F'(U_\alpha)
  + E_\eps^{(2)},
  \label{eq:Wtilde-equation}
\end{equation}
}
with the same remainder \eqref{eq:E2-crude}.

\medskip
\noindent{ 
\emph{Step 2: a refined weighted estimate.}
}
We claim that
\begin{equation}
  |\widetilde W(z)|
  \le
  Cs_\eps\bigl(|\mathbf{b}^{\rm eff}|+s_\eps\bigr)(1+|z|)^\sigma, 
  \qquad
  |z|\le\frac{3r}{4s_\eps}.
  \label{eq:Wtilde-refined}
\end{equation}
To prove this, apply the full-disk normalization and Green argument
of Lemma~\ref{lem:preliminary-pointwise} to
$\widetilde W=W-\Xi_1$, with comparison scale
$\mathfrak n_\eps=s_\eps(|\mathbf{b}^{\rm eff}|+s_\eps)$. The source in
\eqref{eq:Wtilde-equation} divided by $\mathfrak n_\eps$ is bounded by
$C(1+|z|)^{2-m_*}$, and the boundary oscillation is
$O(\etae^2)+O(s_\eps^2\de^2)=O(\mathfrak n_\eps)$.
Consequently the normalized forcing and boundary terms vanish in a
contradiction sequence. Local limits vanish by
Lemma~\ref{thm:linearized-classification}, and the Green tails are
uniformly small since $m_*>4+\sigma$. This proves
\eqref{eq:Wtilde-refined}; its weighted bound holds on the full disk
$D_\eps$.

We next claim that
\begin{equation}
  |\mathbf{b}^{\rm eff}|\le Cs_\eps.
  \label{eq:beff-O-s}
\end{equation}
Otherwise, along a subsequence,
$|\mathbf{b}^{\rm eff}|/s_\eps\to\infty$,
$\mathbf{b}^{\rm eff}/|\mathbf{b}^{\rm eff}|\to\mathbf e_*$, $|\mathbf e_*|=1$, and
$\alpha_\eps\to\alpha_*$. The functions
$\widetilde W/(s_\eps|\mathbf{b}^{\rm eff}|)$ converge locally to $\Psi$ with
\begin{equation}\label{eq:Psi-effective-limit}
 L_{\alpha_*}\Psi=-(\mathbf e_*\cdot z)F'(U_{\alpha_*}),\qquad
 \Psi(0)=0,\quad\nabla\Psi(0)=0,\quad
 |\Psi(z)|\le C(1+|z|)^\sigma.
\end{equation}
Rotate so that $\mathbf e_*=(1,0)$. The cosine first mode $\Psi_1$ satisfies
the forced $k=1$ equation, so $\Psi_1(\varrho)+\varrho$ solves its homogeneous
equation. Regularity at zero implies
$\Psi_1(\varrho)=-\varrho+cU_{\alpha_*}'(\varrho)=-\varrho+O(\varrho^{-1})$, contradicting
$\sigma<1$. This proves \eqref{eq:beff-O-s}.

Combining \eqref{eq:beff-O-s} with
\eqref{eq:Wtilde-refined}, we obtain
\begin{equation}
  |\widetilde W(z)|
  \le
  C\eta_\eps^2(1+|z|)^\sigma,
  \qquad
  |z|\le\frac{3r}{4s_\eps}.
  \label{eq:Wtilde-O2-weighted}
\end{equation}
Moreover, Lemma~\ref{lem:weighted-green}, together with the standard
interior gradient estimate for the harmonic boundary part, gives
\begin{equation}
  |\nabla\widetilde W(z)|
  \le
  C\eta_\eps^2,
  \qquad
  |z|\le\frac{3r}{4s_\eps}.
  \label{eq:grad-Wtilde-O2}
\end{equation}

\medskip
\noindent
{ 
\emph{Step 3: sharp first Fourier estimate.}
}
For $\Pi_1f(\varrho)=(2\pi)^{-1}\int_0^{2\pi}
f(\varrho \boldsymbol\omega(\theta))\boldsymbol\omega(\theta)\,d\theta$, set
\begin{equation}\label{eq:P1-definition}
 Z_{\rm eff}(z)=-s_\eps \mathbf{b}^{\rm eff}\cdot z
 -\frac{2s_\eps}{\kappa}\mathbf{b}^{\rm eff}\cdot\nabla U_\alpha(z),
 \qquad \mathcal R=\widetilde W-Z_{\rm eff}.
\end{equation}
Then $L_\alpha\mathcal R=E_\eps^{(2)}$ and
$\mathcal R(0)=\nabla\mathcal R(0)=0$.
To make the first-mode cancellation explicit, put
\[
 \begin{aligned}
 a^{(1)}(z)&=\Xi_1(z)+s_\eps\mathbf b\cdot z,
 &\Gamma^{(1)}(z)&=s_\eps\mathbf g\cdot z,\\
 \ell^{(2)}(z)&=\frac{s_\eps^2}{2}z^T\mathsf Bz,
 &\Gamma^{(2)}(z)&=\frac{s_\eps^2}{2}z^T\mathsf Cz.
 \end{aligned}
\]
where $\mathsf B,\mathsf C$ are defined in \eqref{eq:B-tracefree}.
The quadratic coefficient terms and the quadratic products of first
modes contribute $-\mathcal Q_2$ to $E_\eps^{(2)}$, where
\[
 \begin{aligned}
 \mathcal Q_2={}&F'(U)\ell^{(2)}+F(U)\Gamma^{(2)}
  +\frac12F''(U)(a^{(1)})^2 +F'(U)\Gamma^{(1)}a^{(1)}
  +\frac12F(U)(\Gamma^{(1)})^2.
 \end{aligned}
\]
At each fixed radius, $a^{(1)}$ and $\Gamma^{(1)}$ are pure first
Fourier modes, whereas $\ell^{(2)}$ and $\Gamma^{(2)}$ contain only
modes $0$ and $2$. Products of two first modes also contain only
modes $0$ and $2$, so $\Pi_1\mathcal Q_2=0$ exactly.
On $|z|\le3r/(4s_\eps)$, the remaining terms are estimated using
\[
 \begin{aligned}
 |a^{(1)}|+|\Gamma^{(1)}|\le Cs_\eps(1+|z|),\ \ \
 |\ell^{(2)}|+|\Gamma^{(2)}|\le Cs_\eps^2(1+|z|)^2,\ \ \
 |\widetilde W|\le Cs_\eps^2(1+|z|)^\sigma.
 \end{aligned}
\]
The coefficient Taylor remainders and cubic products give the power
$(1+|z|)^{3-m_*}$ below. The remaining mixed products with
$\widetilde W$ and its square give the powers
$(1+|z|)^{2+\sigma-m_*}$, $(1+|z|)^{1+\sigma-m_*}$, and
$(1+|z|)^{2\sigma-m_*}$; higher products are absorbed using
$s_\eps(1+|z|)\le C$ and $\sigma<1$.
Together with \eqref{eq:E2-crude} on the outer annulus, this yields
\begin{align}
 |\Pi_1E_\eps^{(2)}(\varrho)|
 &\le Cs_\eps^3\big[(1+\varrho)^{3-m_*}+(1+\varrho)^{2+\sigma-m_*}
       \notag\\[-1mm]
 &\hspace{32mm} +(1+\varrho)^{1+\sigma-m_*}+(1+\varrho)^{2\sigma-m_*}\big],
 &&\varrho\le\frac{3r}{4s_\eps},\label{eq:P1-error-after-refinement}\\
 |\Pi_1E_\eps^{(2)}(\varrho)|
 &\le Cs_\eps^2(1+\varrho)^{2-m_*},
 &&\frac{3r}{4s_\eps}\le \varrho\le\frac{r_\eps}{s_\eps},
 \label{eq:P1-error-outer}
\end{align}
where $r_\eps=r-|\mathbf{d}|$. Let $h_1=U_\alpha'$ and choose $h_2$ so that
$\mathcal W(h_1,h_2)=1/\varrho$; uniformly on $I$,
$|h_1|\le C/(1+\varrho)$ and $|h_2|\le C(\varrho+\varrho^{-1})$.
Variation of parameters gives
\begin{equation}\label{eq:P1-R-variation}
 \Pi_1\mathcal R(\varrho)
 =-h_1(\varrho)\int_0^\varrho h_2(\tau)\Pi_1E_\eps^{(2)}(\tau)\tau\,d\tau
 +h_2(\varrho)\int_0^\varrho h_1(\tau)\Pi_1E_\eps^{(2)}(\tau)\tau\,d\tau.
\end{equation}
Since $m_*>4$ and $\sigma<m_*-4$, these integrals satisfy
\begin{align}
 \int_0^{r_\eps/s_\eps}|h_1(\tau)|
       |\Pi_1E_\eps^{(2)}(\tau)|\tau\,d\tau
 &\le Cs_\eps^3,\label{eq:P1-growing-moment}\\
 |h_1(\varrho)|\int_0^\varrho|h_2(\tau)|
       |\Pi_1E_\eps^{(2)}(\tau)|\tau\,d\tau
 &\le Cs_\eps^2
 \quad(1\le \varrho\le r_\eps/s_\eps).
 \label{eq:P1-decaying-part}
\end{align}
The outer contribution in the first estimate is
$O(s_\eps^{m_*-1})=o(s_\eps^3)$. The growing term is also
$O(s_\eps^2)$ because $h_2(\varrho)=O(\varrho)$, and hence
\begin{equation}\label{eq:first-mode-remainder}
 |\Pi_1\mathcal R(\varrho)|\le Cs_\eps^2
 \qquad(1\le \varrho\le r_\eps/s_\eps).
\end{equation}
At $\varrho=r_\eps/s_\eps$, use $\mathbf{b}^{\rm eff}=O(s_\eps)$,
$U_\alpha'(\varrho)=O(s_\eps)$, and
$\Pi_1\Xi_1(\varrho)=O(s_\eps^2\de^2)$ to obtain
\begin{align}
 \Pi_1\widehat v(r_\eps/s_\eps)
 &=-\frac{r_\eps}{2}\mathbf{b}^{\rm eff}+O(s_\eps^2),
 \qquad 
 \Pi_1\widehat v(r_\eps/s_\eps)
 =-\frac m{2r}\mathbf{d}+O(\etae^2)=O(\etae^2).
 \label{eq:P1-green-boundary}
\end{align}
The second line is Lemma~\ref{lem:first-moment}. Comparing the two,
with $r_\eps=r+O(\etae^2)$ and $s_\eps\asymp\etae$, proves
$|\mathbf{b}^{\rm eff}|\le C\etae^2$, namely assertion~(ii).

\medskip\noindent
\emph{Step 4: the center shift.}
By \eqref{eq:Wtilde-O2-weighted}--\eqref{eq:grad-Wtilde-O2},
$\|W-\Xi_1\|_{C^1(B_R)}\le C_R\etae^2$ for every fixed $R$.
At $\zeta_\eps=-\mathbf d/s_\eps=O(\etae)$, the expansion
$Q_\alpha(\varrho)=O(\varrho^3)$ gives
$|\nabla\Xi_1(\zeta_\eps)|\le C\etae^3\de^2$, while
$\nabla U_\alpha(\zeta_\eps)=-\kappa\zeta_\eps/2+O(|\zeta_\eps|^3)$.
Insert these estimates and
$\nabla\mathfrak h(q)=\nabla\mathfrak h(q^*)+O(\etae^2)$
into \eqref{eq:max-condition-scaled} and multiply by $s_\eps\asymp\etae$:
\(
 \mathbf d=-\frac{2s_\eps^2}{\kappa}\nabla\mathfrak h(q^*)
             +O(\etae^3+\etae^4\de^2).
\)
This proves assertion~\textup{(iii)}.

\medskip
\noindent
\emph{Step 5: the second Fourier mode.}
For $\mathsf B,\mathsf C$ from \eqref{eq:B-tracefree}, Taylor expansion yields
\begin{align}
 \ell(z)=\log A(z)
 &=s_\eps \mathbf{b}\cdot z+\frac{s_\eps^2}{2}z^T\mathsf B z+O(s_\eps^3|z|^3),
 \label{eq:logA-second-correct}\\
 \Gamma(z)
 &=s_\eps \mathbf{g}\cdot z+\frac{s_\eps^2}{2}z^T\mathsf C z
   +O(s_\eps^3\de^3|z|^3).
 \label{eq:logGamma-second-correct}
\end{align}
The decomposition $z^T\mathsf S z=\tfrac12(\tr\mathsf S)|z|^2+z^T\mathsf S^\circ z$
separates modes $0$ and $2$, giving
\begin{align}
 e^\Gamma F(U+W+\ell)
 ={}&F(U)+F'(U)W+s_\eps[F'(U)\mathbf{b}+F(U)\mathbf{g}]\cdot z
 \notag\\
 &+\frac{s_\eps^2}{4}[F'(U)\tr\mathsf B+F(U)\tr\mathsf C]|z|^2
 \notag\\
 &+\frac{s_\eps^2}{2}[F'(U)z^T\mathsf B^\circ z+F(U)z^T\mathsf C^\circ z]
   +\mathcal R_\eps^{(2)}+\mathcal R_\eps^{(3)},
 \label{eq:full-second-order-expansion}
\end{align}
where $\mathbf{b}=O(\etae^2+\de^2)$, $W=\widetilde W+\Xi_1$, and the
established weighted bounds give
\begin{align}
 |\mathcal R_\eps^{(2)}|
 &\le C\big[\etae^4(1+|z|)^{-m_*+2\sigma}
       +\etae^4(1+|z|)^{-m_*+4}
 +(\etae^6+\etae^2\de^4)(1+|z|)^{-m_*+2}\big],
 \label{R1}\\
 |\mathcal R_\eps^{(3)}|
 &\le C\etae^3(1+|z|)^{-m_*+3}.
 \label{R2}
\end{align}
For either angular component of \eqref{eq:Psi2-equation}, let $f(\varrho)$
be its radial right-hand side. Let $h_0\sim\varrho^2$ at zero and
$h_\infty\sim\varrho^{-2}$ at infinity solve the homogeneous $k=2$ equation.
Their constant Wronskian factor
$\omega=\varrho(h_0h_\infty'-h_0'h_\infty)$ is nonzero by
Lemma~\ref{thm:linearized-classification}. The regular solution with
no growing branch is
\[
 \psi(\varrho)=\frac1\omega\left[
 h_\infty(\varrho)\int_0^\varrho t h_0(t)f(t)\,dt
 +h_0(\varrho)\int_\varrho^\infty t h_\infty(t)f(t)\,dt\right].
\]
Indeed, $f(\varrho)=O(\varrho^2)$ at zero and
$f(\varrho)=O(\varrho^{2-m_*})$ at infinity;
the formula gives $\psi=O(\varrho^2)$ at zero and bounded
$|\psi|+\varrho|\psi'|$ at infinity since $m_*>4$.
The bounds are uniform for $\alpha\in I$ by compactness and
$\omega\ne0$; the same kernel classification gives uniqueness.
Combining the two angular components defines $\Psi_2$ and gives
\begin{equation}\label{eq:Psi2-bound}
 \Psi_2(0)=0,\quad\nabla\Psi_2(0)=0,\qquad
 |\Psi_2(z)|+(1+|z|)|\nabla\Psi_2(z)|\le C.
\end{equation}
For \(
 \mathcal T=W-\Xi_1-s_\eps^2\Psi_2,\)
we obtain
\begin{align}
 L_\alpha\mathcal T
 ={}&-\frac{s_\eps^2}{4}[F'(U)\tr\mathsf B+F(U)\tr\mathsf C]|z|^2
 \notag\\
 &-s_\eps(\mathbf{b}^{\rm eff}\cdot z)F'(U)
   -\mathcal R_\eps^{(2)}-\mathcal R_\eps^{(3)},
 \qquad\mathcal T(0)=\nabla\mathcal T(0)=0.
 \label{eq:W-second-equation-correct}
\end{align}
Since $|\mathbf{b}^{\rm eff}|=O(\etae^2)$,
$|\tr\mathsf B|+|\tr\mathsf C|=O(\de^2)$, and $|z|\le C/\etae$,
\begin{equation}\label{eq:second-mode-equation-bound}
 |L_\alpha\mathcal T|
 \le C[\etae^3(1+|z|)^{-m_*+3}
       +\etae^2\de^2(1+|z|)^{-m_*+2}].
\end{equation}
Set $\tau_\eps=\etae^{2+\sigma}+\etae^2\de^2$ and choose
$0<\rho<m_*-4-\sigma$. The normalization of the first source term in
\eqref{eq:second-mode-equation-bound} is justified, uniformly for
$|z|\le C/\etae$, by
\[
 \etae^{1-\sigma}(1+|z|)^{5+\rho-m_*}
 \le C\bigl(\etae^{1-\sigma}
             +\etae^{m_*-4-\sigma-\rho}\bigr)\le C.
\]
Indeed, use $1+|z|\ge1$ when $5+\rho-m_*\le0$ and
$1+|z|\le C/\etae$ otherwise. Both powers of $\etae$ are positive,
since $\sigma<1$ and $m_*-4-\sigma-\rho>0$.
For the second source term,
\(\frac{\etae^2\de^2}{\tau_\eps} (1+|z|)^{4+\rho-m_*}\le1,\)
because $\tau_\eps\ge\etae^2\de^2$ and $4+\rho-m_*<0$.
Consequently,
$\tau_\eps^{-1}|L_\alpha\mathcal T|
\le C(1+|z|)^{-2-\rho}$, without any relation between $\etae$ and $\de$.
The harmonic extension $h_\eps$ of $\mathcal T|_{\partial D_\eps}$
has oscillation $O(\etae^2)$. Recall that
$D_\eps=B_{R_\eps}(\zeta_\eps)$, with
$R_\eps\asymp\etae^{-1}$ and $\zeta_\eps=O(\etae)$.
On $|z|\le R_\eps/2$, the interior gradient estimate gives
$|h_\eps(z)-h_\eps(0)|\le C\etae^3|z|$.
On the remaining part of $D_\eps$, the same inequality follows from
the oscillation bound and $|z|\ge R_\eps/2$.
Thus, using $\sigma<1$ and $\tau_\eps\ge\etae^{2+\sigma}$,
\[
 |h_\eps(z)-h_\eps(0)|\le C\etae^3|z|
 \le C\tau_\eps(1+|z|)^\sigma\qquad(z\in D_\eps).
\]
The full-disk Green and normalized-kernel argument of
Lemma~\ref{lem:preliminary-pointwise} now gives
$|\mathcal T|\le C\tau_\eps(1+|z|)^\sigma$.
This proves assertion~(iv).

\medskip\noindent
\emph{Step 6: the local mass and uniform estimate.}
By \eqref{eq:correct-source-change-variable},
$2\pi\widetilde m_{t,\eps}=\int_{D_\eps}\mathcal F_\eps$.
In \eqref{eq:full-second-order-expansion}, the first and trace-free
second modes integrate to zero on centered disks. The displacement
$\zeta_\eps=O(\etae)$ of $D_\eps$ contributes only a decaying tail;
indeed, the omitted radial mass is $O(\etae^{m_*-2})$.
The trace terms are $O(\etae^2\de^2)$, and the largest remainder
integrals are
\[
 \etae^3\int_0^{C/\etae}(1+\varrho)^{3-m_*}\varrho\,d\varrho
 +\etae^4\int_0^{C/\etae}(1+\varrho)^{4-m_*}\varrho\,d\varrho
 =O(\etae^{2+\sigma}).
\]
At $m_*=5$ or $6$ the additional logarithm is absorbed by the strict
choice \eqref{eq:sigma-choice}. Moreover,
$W=\mathcal T+\Xi_1+s_\eps^2\Psi_2$, so the same angular cancellation
and the weighted bound on $\mathcal T$ give
\[
 2\pi\widetilde m_{t,\eps}
 =2\pi m_{t,\eps}+\int_{D_\eps}F'(U_\alpha)W\,dz+O(\tau_\eps)
 =2\pi m_{t,\eps}+O(\tau_\eps).
\]
This proves assertion~(v).
Finally, on $|z|\le r/(2s_\eps)$, since $\sigma<1$,
\[
 \tau_\eps(1+|z|)^\sigma
 \le C(\etae^2+\etae^{2-\sigma}\de^2)
 \le C(\etae^2+\etae\de^2).
\]
Hence $|W|\le|\mathcal T|+|\Xi_1|+s_\eps^2|\Psi_2|
\le C(\etae^2+\etae\de^2)$, and interior elliptic estimates give
the same bound for $\nabla W$.
If $\varphi_i$ is constant, $\Xi_1$ and the metric trace vanish;
$\log H$ is already harmonic by its definition. This proves assertion~(i) and
completes the theorem.
\end{proof}

\end{document}
